% DEPRECATED：禁止分章导入。请使用同目录的论文.tex 单一总稿。
\PackageError{cumcm-paper}{Deprecated split chapter asset}{Use the single master paper.tex}
\endinput
\section{模型的建立与求解}

% 多问共享同一坐标系、状态方程或基础关系时，可先增加一个信息明确的共享模型二级标题。

\subsection{问题一}

% 从当前题目的对象、条件、关系或实际计算动作起句；不强制写路线段。
% 不强制每问都写“输入—转化—输出”小摘要；审计、原稿、重构、口径、闭合和证据边界不得进入正文。
% 每问在进入下一问前完成“模型关系—具体求解—结果解释—独立验证—风险边界—直接结论”的职责闭环；
% 这些是内容责任，不是必须逐项建立的标题。

\subsubsection{受力分析}

% 定义对象、变量、目标与约束，给出本题关键关系、边界条件和必要推导。

\subsubsection{数值求解}

% 写明输入、变换或离散、参数求取、迭代更新、约束处理、停止条件和报告精度。
% 国一模式：若本问包含三阶段以上、判断/循环/回退/早停、粗搜—细化、递归再生或状态转移，
% 必须在本问分析或求解段附近插入分问题流程图；图中变量、阈值、更新和停止条件须与正文及脚本一致。
% 每问按待证明命题在文字、公式、表格和图之间选择载体；图数不设上下限，使用后须与正文形成合理交互，不套固定前后分析模板。

\subsection{问题二}

% 成稿时直接从问题二的对象、条件或实际计算动作起句，再按实际内容决定是否设置三级标题。

% 按实际分问增删。单问较短时可不设三级标题，直接连续写模型与求解过程。
